\documentclass[a4paper,12pt,twoside]{article}

% Paquetes básicos
\usepackage[utf8]{inputenc}
\usepackage[spanish]{babel}
\usepackage{graphicx}
\usepackage{geometry}
\usepackage{hyperref}
\usepackage{fancyhdr}
\usepackage{caption}

% Márgenes y geometría
\geometry{a4paper, margin=2.5cm}

% Configuración del estilo de página
\pagestyle{fancy}
\fancyhf{}
\fancyhead[LE,RO]{\thepage}
\fancyhead[RE,LO]{Memoria de la Práctica}

% Redefinir cleardoublepage para asegurar página en blanco sin numeración
\makeatletter
\def\cleardoublepage{\clearpage\if@twoside \ifodd\c@page\else
    \hbox{}
    \vspace*{\fill}
    \begin{center}
       % Opcional: \textit{Página intencionadamente en blanco}
    \end{center}
    \vspace{\fill}
    \thispagestyle{empty}
    \newpage
    \if@twocolumn\hbox{}\newpage\fi\fi\fi}
\makeatother

\begin{document}

% ==============================================
% PÁGINA 1: PORTADA
% ==============================================
\begin{titlepage}
    \thispagestyle{empty} % Ocultar numeración en portada
    \begin{center}
        \vspace*{2cm}
        
        % Logo URJC
        % \includegraphics[width=0.6\textwidth]{logo_urjc.png} \\ % Sustituir por la ruta real del logo
        
        \vspace{2cm}
        
        \Huge \textbf{Práctica 6} \\
        \vspace{0.5cm}
        \Large \textbf{Asignatura: Inteligencia Artificial}
        
        \vspace{3cm}
        
        \Large \textbf{Autores:} \\
        \vspace{0.5cm}
        Nombre del Estudiante 1 - DNI: XXXXXXXX-X \\
        Nombre del Estudiante 2 - DNI: XXXXXXXX-X \\ % Añadir o quitar autores según corresponda
        
        \vspace{3cm}
        
        \Large \textbf{Universidad Rey Juan Carlos} \\
        \large Grado en Ingeniería Informática \\
        \large Curso 2025-26
        
    \end{center}
\end{titlepage}
% ==============================================

% Reiniciar contador de páginas internamente DESPUÉS de la portada
% Así la página en blanco siguiente es la 0, y el índice la 1
\setcounter{page}{0}

% ==============================================
% PÁGINA 2: EN BLANCO (Manejada por twoside y cleardoublepage)
% ==============================================
\cleardoublepage

% ==============================================
% PÁGINA 3: ÍNDICE
% ==============================================
\thispagestyle{empty} % Ocultar numeración en índice
\tableofcontents

% Si hay 5 o más figuras, se debe descomentar esto:
% \listoffigures

\cleardoublepage

% ==============================================
% SECCIONES OBLIGATORIAS (Cada una comienza en impar)
% ==============================================

% La numeración visible empieza a partir de la Introducción
\pagestyle{fancy}

\section{Introducción}
% Normalmente de media página a una página completa.
Aquí comienza la introducción de la práctica. Explicar brevemente el problema a resolver y los objetivos propuestos.

\cleardoublepage

\section{Contexto}
Descripción del entorno de trabajo, herramientas utilizadas y bases teóricas de la práctica.

\cleardoublepage

\section{Estado del Arte}
Investigación y análisis de aproximaciones existentes, posibles soluciones y alternativas tecnológicas.

\cleardoublepage

\section{Propuesta}
Explicación detallada de la arquitectura, diseño e implementación realizada para resolver el problema.
Se puede dividir en subsecciones si es necesario.

\cleardoublepage

\section{Conclusiones}
% Normalmente de media página a una página completa.
Resumen de los resultados obtenidos, dificultades encontradas y posibles mejoras futuras.

\cleardoublepage

% ==============================================
% REFERENCIAS
% ==============================================
\begin{thebibliography}{9}
\bibitem{ref1} Autor. \textit{Título}. Editorial, Año.
\end{thebibliography}

\end{document}
